\documentclass[11pt]{article}
\usepackage[a4paper,margin=27mm]{geometry}
\usepackage{amsmath,amssymb,mathtools}
\usepackage{fontspec,unicode-math}
\setmainfont{Latin Modern Roman}
\setmathfont{Latin Modern Math}
\usepackage{xurl,hyperref}
\hypersetup{colorlinks=true,urlcolor=blue,linkcolor=blue}
\newcommand{\sref}[2]{\hyperlink{source:#1:#2}{[S#2]}}
\newcommand{\cataloguescope}[1]{\paragraph{Mathematical question.} #1}
\begin{document}
\section{Conjecture 4.4 — Amoroso-David}\label{statement}

\subsection{Definitions and mathematical statement}

\paragraph{Shared notation.}$\mathbb N=\{1,2,\ldots\}$, and $\mathbb Z,\mathbb Q,\mathbb R,\mathbb C$ denote the integers, rationals, reals and complex numbers. Any other conventions are specified locally.


Let $\alpha=(\alpha_1,\ldots,\alpha_n)\in(\overline{\mathbb Q}^{\times})^n$. Multiplicative independence means $\prod_i\alpha_i^{a_i}=1$ with $a_i\in\mathbb Z$ implies all $a_i=0$. Define $\omega(\alpha)=\min\{\deg P:0\ne P\in\mathbb Q[X_1,\ldots,X_n],\ P(\alpha)=0\}$. For any number field $K$ containing the coordinates, the absolute logarithmic projective height is $h(1:\alpha_1:\cdots:\alpha_n)=[K:\mathbb Q]^{-1}\sum_v[K_v:\mathbb Q_v]\log\max\{1,|\alpha_1|_v,\ldots,|\alpha_n|_v\}$, with the usual normalized absolute values.
\paragraph{Statement recorded in the supplied source.}
For every $n\ge1$, does there exist $c(n)>0$ such that every multiplicatively independent algebraic tuple $\alpha$ satisfies \[h(1:\alpha_1:\cdots:\alpha_n)\ge c(n)/\omega(\alpha)?\] \sref{statement}{2}
\subsection{Short English statement}
Is the projective height of a multiplicatively independent algebraic tuple uniformly bounded below by a constant divided by the least degree of a rational polynomial vanishing at that tuple?
\subsection{Sources}
\begin{description}
\item[\hypertarget{source:statement:1}{S1}] ULAM.AI and the UnsolvedMath Contributors, UnsolvedMath: A Curated Collection of Open Mathematics Problems, catalogue entry. CC BY 4.0. \url{https://huggingface.co/datasets/ulamai/UnsolvedMath}.
\item[\hypertarget{source:statement:2}{S2}] Michel Waldschmidt, Open Diophantine Problems, arXiv:math/0312440v2 (24 January 2004), Conjecture 4.4 and the definitions at the start of Section 4. \url{https://arxiv.org/abs/math/0312440}.
\end{description}
\label{end:statement}
\end{document}
